\subsection{B. Die einfache Darstellung vermischt zwei
Näherungszustände}\label{b.-die-einfache-darstellung-vermischt-zwei-nuxe4herungszustuxe4nde}

Die Werte um -1.138476Δ, mittlere Bosonzahl 1.0177445, Überlappung mit F
von 0.3498904 und Bosonzahl-Shannonwert 1.8712436 Bit gehören zum
größeren siebendimensionalen Ritz-Zustand. Die im selben Text genannten
Gewichte 0.973684 und 0.026316 gehören dagegen ausdrücklich zu der
älteren fünfdimensionalen Basis v0 bis v3 plus w2. Deren Energie liegt
um -1.0942319Δ und ihre mittlere Bosonzahl um 0.8421159.

Die unterschiedlichen Werte wurden aus den beiden kleinen Matrizen
unabhängig nachgerechnet. Für den \textbf{größeren selben Ritz-Zustand}
folgen stattdessen die gesamten Entnahme-/Additionsnormen ungefähr
0.9681955 und 0.0318045. Keine dieser Näherungszahlen ist das Gewicht
des isolierten Pols auf dem wahren stationären Grundzustand. Auch die
angegebenen Entnahme-/Additionsenergiemittel aus der kleineren Basis
sind keine exakten Polenergien.

Eine zusätzliche CAR-Summenregelrechnung bestimmt die ersten geladenen
Momente aus genau denselben kleinen Matrizen, ohne Millionen
Besetzungsamplituden neu zu konstruieren. Sie reproduziert die
gespeicherten Mittelwerte der Fünferbasis und ergibt für den größeren
Siebener-Ritz-Zustand stattdessen mittlere Entnahmekosten von etwa
0.03479766990Δ und Additionskosten von etwa 1.05931330814Δ. Das sind
korrigierte numerische Variationszustandswerte, keine Ω-Polenergien. Die
allgemeinen Identitäten und die reproduzierbare Auswertung stehen im
ausführlichen Nachtrag A.

Die Aussage „unser schwächerer Sektorvergleich schließt die Lücke
nicht`` ist korrekt. Sie widerlegt aber keinen stärkeren vorherigen
Grundzustandssatz für denselben Modellvertrag. Das eingefrorene fremde
Gesamt-Replaymanifest stand noch auf RUNNING. Deshalb wird hier nur
unser abgeschlossener isolierter Norm- und Ritz-Replay behauptet, nicht
der Abschluss der fremden Gesamtrunde.

Der Text nennt eine fast geschlossene Leiter. Der bereits bewiesene
Rücklaufrest bleibt jedoch ungleich null. 33 mögliche Bosonzahlstufen
bedeuten nicht 33 Vektoren für den gesamten Prozess. Auch alle 33
skalaren Normen würden zunächst einen bestimmten Variationsraum
beschreiben, nicht automatisch die volle Antwort.

\subsection{C. Ein neuer entscheidender Test am
Feldprojektor}\label{c.-ein-neuer-entscheidender-test-am-feldprojektor}

Der Spin-1/2-Projektor im zweiten Text trägt unter einer bestimmten
Händigkeit. Für

\[\adjustbox{max width=\linewidth}{$\displaystyle 
(1,0)\otimes(1/2,0)=(3/2,0)\oplus(1/2,0)
$}\]

wurde eine explizite Epsilon-Auslese C samt Wiedereinsetzung R
konstruiert. Alle gleichhändigen sl(2)-Intertwining-Gleichungen gelten,
CR=I2 und RC=P1/2. Die Gram-Matrix der sechs verwendeten unnormierten
Koordinaten ist diag(1,1,2,2,1,1); diese Metrik darf nicht durch die
Einheitsmatrix ersetzt werden.

Wird der zweite Faktor jedoch als \textbf{wörtlich adjungiertes linkes
Weylfeld} gelesen, hat er entgegengesetzte Händigkeit. Dann gilt

\[\adjustbox{max width=\linewidth}{$\displaystyle 
(1,0)\otimes(0,1/2)=(1,1/2),
$}\]

ein irreduzibler Lorentzraum. Unter Drehungen gibt es weiterhin Spin 3/2
und Spin 1/2; der Rotationsprojektor ist aber nicht boostinvariant. Der
neue exakte Prüfer findet

\[\adjustbox{max width=\linewidth}{$\displaystyle 
[P_{1/2},J_z]=0,\qquad
\operatorname{rank}[P_{1/2},K_z]=4.
$}\]

Das ist ein konkreter Test, weshalb „das Feldwörterbuch ist fertig`` zu
weit geht. Es ist \textbf{kein pauschaler Ausschluss des ursprünglichen
Modenkomposits}: Ein Fock-Erzeuger f† ist nicht automatisch ein
adjungiertes lokales Weylfeld. Eine linke Weylfeldentwicklung enthält
bereits Erzeugungs- und Vernichtungsterme mit passenden
Spinorwellenfunktionen. Eine unabhängige oder korrekt ladungskonjugierte
gleichhändige Konstruktion bleibt möglich; sie benötigt den
tatsächlichen Adapter mit Ladung, CAR, W und Dynamik. Die Konventionen
entsprechen \href{https://arxiv.org/abs/0812.1594}{Dreiner, Haber und
Martin, Abschnitte 2 und 3}.

Der nächste Test ist damit konkret: denselben geladenen Operator samt
gepunkteten/ungepunkteten Indizes abbilden, sowohl Drehungen als auch
Boosts verschränken, dann sein wirkliches Matrixelement am nativen
Grundzustand bestimmen. Ein Projektorrang von zwei aus sechs ist kein
zustandsunabhängiges Spektralgewicht von einem Drittel. Positive
Zustände mit Gewicht null und eins sind explizit geprüft.

\subsection{D. Weitere
Übernahmegrenzen}\label{d.-weitere-uxfcbernahmegrenzen}

\begin{itemize}
\tightlist
\item
  Die Clock ist nach der späteren konkreten Konstruktion ein Element der
  inneren Spin(10)-Darstellung. Die ältere Formulierung „äußerer
  Handgriff`` wird nicht als neue Herkunftstatsache übernommen.
\item
  Die große Kommutantdimension zählt nicht direkt physische
  Freiheitsgrade. Ein skalarer Kommutant bedeutet nicht automatisch
  beliebige dynamische Kontrolle. Als exakter Gegenzeuge haben
  Spin-1-Kontrollen einen skalaren Kommutanten, aber Lie-Dimension drei
  statt acht.
\item
  Die Viererregel N=0 modulo vier ist eine notwendige Singulettregel.
  Ohne globalen Eindeutigkeitssatz darf sie tiefere entartete
  Nichtsingulett-Grundzustände nicht ausschließen.
\item
  Das skalare Grassmann-Nullresultat ist korrekt in der angegebenen
  lokalen, ableitungsfreien, gleichhändigen Ein-Kopie-Klasse mit
  unverändertem W. Es verbietet nicht jede denkbare skalare Erweiterung
  mit anderen Feldern, Ableitungen oder unabhängigen Komponenten.
\item
  Die angegebenen getrennten Casimire sind wertvolle minimale Sonden;
  nicht jede sinnvolle Kontrolle erfordert sofort die ganze Gruppe.
  Schon ein einzeln hinreichend trennender Casimir kann die betrachteten
  dunklen Typen unterscheiden.
\end{itemize}

\subsection{E. Die dritte Notiz: richtige Richtung, überwiegend bereits
berücksichtigt}\label{e.-die-dritte-notiz-richtige-richtung-uxfcberwiegend-bereits-beruxfccksichtigt}

Die längere fundamentale Reduktion ergänzt den bereits eingefrorenen
früheren Kommentar um die ausführliche Herleitung. Ihr stärkster Satz
trägt: Bei invariantem H, invariantem Ω und irreduzibler innerer
Entnahmefamilie gilt für die \textbf{gesamte} lokale Antwort

\[\adjustbox{max width=\linewidth}{$\displaystyle 
T^\dagger F(H-E_0)T=\gamma_F I_{64},\qquad
C_{rs}(t)=\delta_{rs}c(t).
$}\]

Dieser Satz umfasst Nebenlinien. Beliebig genaues c(t) erzeugt deshalb
keinen Ortsindex in den 64 inneren Labels. Er verbietet weder die neu
geprüften N=2-Interventionen auf anderen präparierten Zuständen noch
Bewegung auf zusätzlichen Multiplizitätsräumen.

Auch die Unterscheidung zwischen dem Kommutanten der inneren Symmetrie
und demjenigen der verfügbaren Operationen ist richtig. Der eine lässt
symmetrieverträgliche Dynamik auf unabhängigen Vorkommen zu, der andere
beschreibt die Reichweite der gewährten Zugriffe. Viele Vorkommen sind
noch keine räumlichen Orte.

Die bekannte Paritätsstruktur einer Bank wird korrekt eingegrenzt: Wenn
Bosonen unverändert bleiben, erhält der verbundene Paargraph nur globale
Fermionparität, keine echte Teilmengenparität. Unabhängig kopierte
Banken können durch ihre gewählte Zusammensetzung neue lokale
Erhaltungsgrößen bekommen. Deshalb ist die Herleitung einer globalen
Kompositionsregel eine echte Alternative zur bloßen Ergänzung eines
Links zwischen fertigen Kopien.

Die Zustandsfrage darf ebenfalls einfacher gestellt werden: Bei festem
Gesamt-N ist μN nur ein skalares Energieglied. Eine fundamentale Theorie
muss keine unbeobachtbare Energiekonvention auswählen. Sie muss aber den
physisch relevanten Zustand oder die Randbedingung begründen. Dazu wird
kein Experimentator außerhalb des Universums gebraucht, der das
Universum aus dem leeren Fockzustand präpariert. Die interne Präparation
eines konkreten Tests bleibt eine davon verschiedene Aufgabe.

Auch das konkrete Gegenbeispiel zur sektorübergreifenden Zustandswahl
ist korrekt. Die beibehaltene Voll-Fock-Casimirschranke ergibt bei
g/Δ=1/20 und μ/Δ=1/50

\[\adjustbox{max width=\linewidth}{$\displaystyle 
H_\mu\ge\left(\mu-\frac{15g^2}{2\Delta}\right)N_f+2\mu N_b
\ge\frac{\Delta}{800}N.
$}\]

Der leere Zustand ist dann eindeutig energetisch bevorzugt. Die
Quellenpins des früheren Zertifikats wurden abgeglichen; sein großer
Lauf wurde nicht wiederholt. Der genaue Koeffizient 1/800 wurde erneut
rational geprüft. In der Interventionsformulierung ist
„nichtverschwindende Differenz`` die passende Bedingung; auch unser
negativer Effekt ist ein Nachweis.

\textbf{Übernahmeentscheidung:} Diese Notiz trägt als Priorisierung.
Eine zusätzliche globale Quell-/Kompositionsregel muss voneinander
unterscheidbare Vorkommen, ihre Kopplung, den Zustand und die internen
Instrumente gemeinsam liefern. Sie ist im Text noch nicht konstruiert.
Der neu gelöste Dreizustandsversuch ist ein kleiner positiver
Referenzvertrag, keine solche wachsende Welt. Den exakten lokalen Pol
weiter zu verfeinern bleibt nützlich, ersetzt aber diese Herkunftsfrage
nicht.

\subsubsection{Automatisch aus dem abgeschlossenen
Replay}\label{automatisch-aus-dem-abgeschlossenen-replay-1}

{\def\LTcaptype{none} % do not increment counter
\begin{longtable}[]{@{}lrr@{}}
\toprule\noalign{}
Eigener Strang & Exakte Bedingungen & Numerische Bedingungen \\
\midrule\noalign{}
\endhead
\bottomrule\noalign{}
\endlastfoot
causal & 662 & 10 \\
common3 & 43 & 0 \\
late\_field & 36 & 0 \\
late\_moments & 2 & 2 \\
late\_norm & 36 & 10 \\
mechanism & 261 & 0 \\
phase\_bridge & 32 & 0 \\
resource\_boundary & 24 & 0 \\
shadows & 5417 & 17 \\
\end{longtable}
}

Je normalem und optimiertem Lauf: 6513 eigene exakte beziehungsweise
faktische und 39 numerische Bedingungen. Die erste Kategorie enthält
Quellenpins und Strukturprüfungen, nicht nur mathematische Identitäten.
Hinzu kommen 17 erhaltene native Konstruktorguards und 1239 erhaltene
Spurnetzwerk-Quellguards, separat 57 Synthese- und 20 historische
Kettenbedingungen.

\begin{center}\rule{0.5\linewidth}{0.5pt}\end{center}

\section{Teilbericht A: ursprünglicher Mechanismus und
Referenzbilanz}\label{teilbericht-a-urspruxfcnglicher-mechanismus-und-referenzbilanz}

\section{Ursprünglicher Mechanismus, Erhaltung und die kleinste
zusätzliche
Ressource}\label{urspruxfcnglicher-mechanismus-erhaltung-und-die-kleinste-zusuxe4tzliche-ressource}

\subsection{Ergebnis}\label{ergebnis-1}

\textbf{Exakt neu:} Am unveränderten Tensor schließt ein
zweidimensionaler G-Singulett-Unterraum mit physischer Ladung N=4 unter
dem gesamten nativen H. Seine Matrix ist

\[\adjustbox{max width=\linewidth}{$\displaystyle 
H_{4,\mathrm{klein}}=
\begin{pmatrix}2\Delta&4g\\4g&\Delta\end{pmatrix}.
$}\]

Der Koeffizient 4 und die Invarianz wurden aus sämtlichen ursprünglichen
Paartermen berechnet. Es handelt sich um einen invarianten Unterraum,
keine abgeschnittene Näherung. Damit liegt ein kleiner ursprünglicher
Mechanismus zwischen zwei verschiedenen Zusammensetzungen desselben
Singuletttyps vor. Seine Vorbereitung aus dem Vakuum ist durch diese
Rechnung nicht geliefert.

\textbf{Exakt ausgeschlossen:} Kein Wort aus X, Nb und den
dokumentierten passiven Focklifts von Clock und innerer Gruppe kann den
Systemsektor N verändern. Das gilt für beliebige Wortlängen. Die
zusätzlichen Terme B+ und R+ aus v1.6.7 können daher in diesem festen
Alphabet nicht durch weitere Kommutatorsuche gefunden werden.

\textbf{Bedingt konstruiert:} Ein zusätzliches Bosonpaarinstrument ist
die einfachste invariant mögliche Erweiterung nach Polynomialgrad.
Zusammen mit einer zweistufigen G-trivialen Referenz entsteht ein exakt
geschlossener Dreizustandsprozess mit voller Spin(10)×SU(4)-Invarianz
und Ngesamt=4. Die Referenz und die gekoppelte Operation sind
ausdrücklich neue Ressourcen. Dieser gemeinsame Prozess implementiert
keinen kostenlosen kohärenten Bosonpaarantrieb auf dem isolierten
System.

\subsection{1. Was die tatsächlichen Quellen
konstruieren}\label{was-die-tatsuxe4chlichen-quellen-konstruieren}

Alle zwölf verwendeten Eingaben sind unter \codeword{inputs/}
unverändert eingefroren; \codeword{inputs\_manifest.json} enthält ihre
ursprünglichen Pfade, Größen und SHA-256. Das maßgebliche Tensorarchiv
besitzt weiter den SHA-256
\codeword{3f00a0892157a691e4ef26920c3702772c7bf76a250fad98ff04a079bb64b763}.

{\def\LTcaptype{none} % do not increment counter
\begin{longtable}[]{@{}
  >{\raggedright\arraybackslash}p{(\linewidth - 4\tabcolsep) * \real{0.3333}}
  >{\raggedright\arraybackslash}p{(\linewidth - 4\tabcolsep) * \real{0.3333}}
  >{\raggedright\arraybackslash}p{(\linewidth - 4\tabcolsep) * \real{0.3333}}@{}}
\toprule\noalign{}
\begin{minipage}[b]{\linewidth}\raggedright
Quellstelle
\end{minipage} & \begin{minipage}[b]{\linewidth}\raggedright
Tatsächlich konstruierter Gegenstand
\end{minipage} & \begin{minipage}[b]{\linewidth}\raggedright
Grenze dieses Anschlusses
\end{minipage} \\
\midrule\noalign{}
\endhead
\bottomrule\noalign{}
\endlastfoot
\codeword{universalraum-singlet-observable-20260915/sources/native\_source.py},
Zeilen 36-67 & Fünf Jordan-Wigner-Hilfsoperatoren, deren
16-dimensionaler gerader Spinorraum, zehn BETA-Matrizen und der Farbkeil
ergeben W & Die Hilfs-CAR auf fünf Bits sind keine Ausführung der 64
physischen Fermionoperatoren \\
Dieselbe Datei, Zeilen 87-99 & BAR und ETA aus entgegengesetztem
Vektorgewicht und komplementärem Farbpaar & Eine invariante
Koeffizientenmatrix ist noch kein Bosonpaarinstrument \\
Dieselbe Datei, Zeilen 101-129 & J und C3 mit JC3=0 sowie JJ†=15I & Ein
weiterer E8-Koeffiziententensor ist keine vollständige
E8-Operatordarstellung auf dem Fockraum \\
\codeword{seam\_state\_derivation\_probe.py}, Helfer ab Zeile 380 und
Hauptkonstruktion 486-627 & Markierte endliche Clockpermutation aus der
ursprünglichen Duad-/Aut-Konstruktion & Keine Auswahl einer
physikalischen Zeiteinheit oder des Hamiltonoperators \\
\codeword{native-operations-ground-response-20260915/common.py}, Zeilen
104-170 und 187-219 & 60 innere Lie-Matrizen, ihr Bosonlift und der
Clocklift & Die konkreten Lifts erhalten Fermion- und Bosonzahl jeweils;
aktive unabhängige Kontrolle bleibt eine weitere Voraussetzung \\
\codeword{operations\_commutant.py}, Zeilen 6-18 & Ausdrücklich
unterschiedene Kontrollstufen: fixes Modell, X/Nb, Clock,
Gruppenmatrizen & Die Stufen klassifizieren gewährte Operatoren, nicht
deren Compilerverfügbarkeit \\
\codeword{compiler-origin-audit-20260913/context\_instrument.py}, Zeilen
1-6 und 109 ff. & Instrumente auf dem ursprünglichen vierdimensionalen
Kontextträger unter Born-/Wiederholbarkeitsannahmen & Kein gegebener
Adapter zu B+, R+ oder der 64f/60b-Bank \\
\codeword{compiler-kernel-foundation-20260914/relational\_kernel.py},
Zeilen 99-121 & Konkrete geordnete Reflexionswörter und ihr interner
Phasenzeuge & Ihre vorhandene C4-Wirkung liefert nicht automatisch einen
ladungsändernden Fockgriff \\
\end{longtable}
}

Die letzten beiden Quellen werden hier zur Typ- und
Verfügbarkeitsprüfung gelesen, nicht erneut als Gesamtprogramme
ausgeführt. Die Quellenprüfung ist auf diese konkret bezeichnete Kette
begrenzt; sie ist keine Behauptung, dass jedes denkbare andere
Compilerabbild oder jede Datei des Repositories ausgeschlossen wurde.
Die verfügbaren MCP-Codegraphtools enthielten keine aufrufbaren
Graphsuchfunktionen; deshalb wurde gezielt lokal gelesen.

Der eigene Replay führt die tatsächliche \codeword{native\_source.py}
bis einschließlich Zeile 220 mit allen 17 dortigen Guards aus. Nur die
nachfolgende optionale Hashbildung des dichten C3-Arrays und der
Druckblock entfallen. Der erhaltene W-Tensor stimmt in jedem Eintrag mit
dem gepinnten Archiv überein. Der ursprüngliche Clockkern wird
einschließlich seiner vier mathematischen S0.1-S0.4-Guards wiederholt;
die übrigen numerischen Seam-Proben laufen nicht. Das ergibt erneut
p=(2,0,1,4,3) und den dokumentierten passiven Focklift.

\subsection{2. Geschlossener Erhaltungssatz für das dokumentierte
Alphabet}\label{geschlossener-erhaltungssatz-fuxfcr-das-dokumentierte-alphabet}

Es gilt auf dem endlichen Besetzungszustandskern

\[\adjustbox{max width=\linewidth}{$\displaystyle 
N=N_f+2N_b,\quad
[N,f_i]=-f_i,\quad [N,b_A]=-2b_A.
$}\]

Mit den unveränderten Definitionen

\[\adjustbox{max width=\linewidth}{$\displaystyle 
P_A=\sum_{i<j}W_{A,ij}f_jf_i,\qquad
T_+=\sum_A b_A^\dagger P_A,\quad T_-=T_+^\dagger,
\quad X=T_++T_-
$}\]

hat jedes der 480 signierten Wechselwirkungsmonome Ladung +2−1−1=0. Nb,
alle Modenbesetzungen und alle passiven Bilineare f†Af beziehungsweise
b†Cb sind ebenfalls neutral. Die dokumentierten inneren Liegeneratoren
werden genau auf diese Weise zweitquantisiert. Die Clock schickt f† nach
GF f† und b† nach GB b†, mischt also keine Erzeuger mit Vernichtern.

Sei A0 die von diesen gewährten Operationen erzeugte *-Algebra. Aus

\[\adjustbox{max width=\linewidth}{$\displaystyle 
[N,AB]=[N,A]B+A[N,B]
$}\]

folgt induktiv

\[\adjustbox{max width=\linewidth}{$\displaystyle 
\boxed{\mathcal A_0\subseteq\{N\}'.}
$}\]

Der Beweis umfasst Produkte, Summen, Adjungierte und Kommutatoren jeder
endlichen Länge. Auf jedem festen N-Sektor sind die betreffenden nativen
Hamiltonoperatoren endlichdimensional; ihre exponentiellen Entwicklungen
und endliche Pulsfolgen erhalten daher ebenfalls diesen Sektor. Es ist
keine numerische Suche nur bis zu einer endlichen Wortlänge.

Für die in v1.6.7 geprüften Kandidaten gilt hingegen

\[\adjustbox{max width=\linewidth}{$\displaystyle 
[N,B_+]=4B_+,\qquad [N,R_+]=4R_+.
$}\]

Ihre Nichtnullwirkung ist direkt sichtbar:

\[\adjustbox{max width=\linewidth}{$\displaystyle 
\|B_+|0\rangle\|^2=30,\qquad
\|R_+|0\rangle\|^2=480.
$}\]

Also gehören B+, R+ und ihre nichttrivialen hermiteschen Quadraturen
nicht zu A0. Ein fixes H bietet nicht mehr Erreichbarkeit als das
großzügigere Alphabet mit X und Nb als unabhängigen Kontrollen. Selbst
alle gewährten inneren Gruppenoperationen beseitigen dieses N-Hindernis
nicht.

